% Direct-authority diplomatic French for printed p.154.
\SourcePage{154}

\LayerTarget{tohoku:def:presheaf-and-sheaf:p2}
\noindent que les $F(U)$ sont des groupes (resp. des anneaux etc.) et que les
$\varphi_{VU}$ sont des homomorphismes, on obtient la notion de préfaisceau
(ou de faisceau) de groupes (resp. d'anneaux etc.) ; plus généralement, on
pourrait définir la notion de préfaisceau ou faisceau à valeurs dans une
catégorie (cf. 1.1) donnée. Les préfaisceaux, ou les faisceaux, sur $X$, à
valeurs dans une catégorie donnée, forment une catégorie, les morphismes étant
définis comme morphismes de systèmes inductifs. Les préfaisceaux ou faisceaux
sur $X$, à valeurs dans une catégorie \textit{additive}, par exemple la
catégorie des groupes abéliens, forment une catégorie additive, et dans le cas
des préfaisceaux et faisceaux de groupes abéliens, même une catégorie
abélienne. (Pour abréger, nous dirons \textit{faisceau abélien} et
\textit{préfaisceau abélien} pour un faisceau ou préfaisceau de groupes
abéliens).

\LayerTarget{tohoku:remark:sheaves-vs-presheaves-abelian}
Mais on fera attention que le foncteur identique, qui à un faisceau abélien
fait correspondre le préfaisceau abélien correspondant, est exact à gauche
mais non pas exact : si on a un homomorphisme de faisceaux
$u:F\longrightarrow G$, son conoyau en tant qu'homomorphisme de préfaisceaux
est le préfaisceau
\[
Q(U)=G(U)/\operatorname{Im}F(U),
\]
qui en général n'est pas un faisceau ; son conoyau en tant qu'homomorphisme de
faisceaux est le faisceau \textit{associé} au préfaisceau $Q$ (voir plus bas).
Nous n'insisterons pas plus sur ces questions, assez bien connues actuellement
(cf. [4] et le livre de Godement [9]).

\LayerTarget{tohoku:def:etale-space-from-presheaf}
Soit $F$ un préfaisceau d'ensembles sur $X$, on pose pour tout $x\in X$ :
\[
F(x)=\varinjlim_{U\ni x}F(U),
\]
la limite inductive étant prise suivant l'ordonné filtrant des vosinages
ouverts $U$ de $x$. Sur l'ensemble $\overline F$, réunion des $F(x)$, on met
la topologie engendrée par l'ensemble des parties de $\overline F$ qui sont de
la forme $A(f)$, où pour tout ouvert $U\subset X$ et tout $f\in F(U)$, on
désigne par $A(f)$ l'ensemble des images canoniques $f(x)$ de $f$ dans les
$F(x)$, $x\in U$. Quand $\overline F$ est muni de cette topologie,
l'application naturelle de $\overline F$ dans $X$ est un homéomorphisme local
(i. e. tout point de $\overline F$ a un voisinage ouvert appliqué
homéomorphiquement sur un ouvert de $X$), nous dirons (en suivant Godement)
que $\overline F$ est un \textit{espace étalé} dans $X$.

\LayerTarget{tohoku:prop:sheaves-etale-spaces-equivalence}
D'ailleurs un espace étalé dans $X$, soit $E$, définit un faisceau
$\mathfrak F(E)$ de façon naturelle, savoir celui qui à un ouvert $U$ associe
l'ensemble des \og section\fg{} continues de $E$ au dessus de $U$ ; de plus
l'espace étalé $\overline{\mathfrak F(E)}$ associé à $\mathfrak F(E)$
s'identifie canoniquement à $E$. Si on part d'un préfaisceau $F$ sur $X$, on
a d'autre part un homomorphisme canonique
\[
F\longrightarrow\mathfrak F(\overline F),
\]
car toute $f\in F(U)$ définit une section continue
$x\longmapsto f(x)$ de $\overline F$ sur $U$ ; cet homomorphisme est un
isomorphisme si et seulement si $F$ est un faisceau. Ces considérations
prouvent : (i) \textit{La notion de faisceau d'ensembles sur $X$ est
équivalente à la notion d'espace étalé dans $X$} (de façon précise, nous avons
défini une équivalance de la catégorie des faisceaux d'ensembles sur $X$, sur
avec la catégorie des espaces étalés dans $X$).

\LayerTarget{tohoku:construction:sheafification:p1}
(ii) \textit{À tout préfaisceau $F$ sur $X$ correspond un faisceau
$\mathfrak F(\overline F)$ et un homomorphisme
$F\longrightarrow\mathfrak F(\overline F)$, qui est un isomorphisme si et
seulement si $F$ est un faisceau.} (D'ailleurs, $\mathfrak F(\overline F)$ est
un

\TohokuContinuation{155}

